
\documentclass{article}
\usepackage{colortbl}
\usepackage{makecell}
\usepackage{multirow}
\usepackage{supertabular}

\begin{document}

\newcounter{utterance}

\centering \large Interaction Transcript for game `wordle\_withclue', experiment `high\_frequency\_words\_clue\_no\_critic', episode 9 with reka{-}flash{-}3.1{-}t0.0.
\vspace{24pt}

{ \footnotesize  \setcounter{utterance}{1}
\setlength{\tabcolsep}{0pt}
\begin{supertabular}{c@{$\;$}|p{.15\linewidth}@{}p{.15\linewidth}p{.15\linewidth}p{.15\linewidth}p{.15\linewidth}p{.15\linewidth}}
    \# & \multicolumn{2}{c}{Player} && \multicolumn{2}{c}{Game Master} \\
    \hline

    \theutterance \stepcounter{utterance}  
    & & & \multicolumn{4}{p{0.6\linewidth}}{
        \cellcolor[rgb]{0.9,0.9,0.9}{
            \makecell[{{p{\linewidth}}}]{
                \texttt{\tiny{[P1$\langle$GM]}}
                \texttt{You are a language wizard who likes to guess words by using the given rules.} \\
\\ 
\texttt{Welcome to Wordle! You have six attempts to guess the target word, a valid English word of five lowercase letters (a{-}z). Please use the tags "explanation:" and "guess:" to provide a concise explanation for each guess.} \\
\\ 
\texttt{To help you make an informed guess, you will receive a clue for the word, such as} \\
\texttt{clue: snowy white.} \\
\\ 
\texttt{Here is an example guess based on the clue:} \\
\texttt{explanation: In the fairy tail Snow White, the girl is killed because she eats a poisoned apple. And the word apple has 5 letters.} \\
\texttt{guess: apple} \\
\\ 
\texttt{After each guess, your answer will be validated, and you will receive feedback indicating which letters are correct (green), which letters are correct but in the wrong position (yellow), and which letters are incorrect (red). This feedback can be useful in determining which letters to include or exclude in your next guess.} \\
\\ 
\texttt{For example, the feedback for "apple" might be:} \\
\texttt{guess\_feedback: a<yellow> p<yellow> p<green> l<yellow> e<red>} \\
\\ 
\texttt{The explanation should contain details about how both the guess\_feedback and clue work together in tandem to formulate a new guess.} \\
\\ 
\texttt{Do not generate any other text.} \\
\\ 
\texttt{Let's begin with your first guess.} \\
\\ 
\texttt{clue: gauge reading} \\
            }
        }
    }
     \\ \\

    \theutterance \stepcounter{utterance}  
    & \multicolumn{4}{p{0.6\linewidth}}{
        \cellcolor[rgb]{0.561,0.737,0.561}{
            \makecell[{{p{\linewidth}}}]{
                \texttt{\tiny{[P1$\rangle$GM]}}
                \texttt{<reasoning>} \\
\texttt{Let's tackle this Wordle with the clue "gauge reading". First, I need to understand what the clue refers to. A gauge reading is a measurement taken from a gauge, which is a device used to measure something. So maybe the target word is related to measurement or a specific type of gauge.} \\
\\ 
\texttt{Since the word is five letters, I need to think of five{-}letter words connected to gauges. Common gauges include things like speedometers, thermometers, barometers, etc. Let's break down the clue. "Gauge" itself is a five{-}letter word, but the clue is "gauge reading", so maybe the word combines parts of these terms or relates to measurement.} \\
\\ 
\texttt{Looking for five{-}letter words. Words like "meter" (but that's six letters), "scale" (five letters). Wait, "scale" could be part of a gauge. Or "gauge" itself. But the clue is "gauge reading", so maybe the word is related to the act of reading a gauge. Words like "read" (but that's four letters). Maybe "reads" but that's five letters. Wait, "reads" is a verb, but maybe the target is a noun. Alternatively, "gauge" plus "read" parts? Not sure.} \\
\\ 
\texttt{Alternatively, think of common five{-}letter words that are associated with measurements. Words like "meter" (but too long), "volt" (voltage), "temp" (temperature), "press" (pressure). But those are longer. Wait, "gauge" itself is five letters. Maybe the answer is "gauge"? Let me check. If the clue is "gauge reading", then "gauge" is part of it. But is "gauge" a valid Wordle word? Yes, it is. But maybe the answer is something else.} \\
\\ 
\texttt{Alternatively, think of "reading" as part of the word. "Reading" is seven letters. Maybe a part of it? "Readi" doesn't make sense. Or "gauge" plus "read" parts. Alternatively, think of "gauge" as a tool, so maybe "gauge" is the answer. Let's test that. If I guess "gauge", but the clue is "gauge reading", maybe the answer is a word related to the process of reading a gauge. Alternatively, maybe "meter" is part of it, but again, too long} \\
            }
        }
    }
    & & \\ \\

    \theutterance \stepcounter{utterance}  
    & & & \multicolumn{2}{p{0.3\linewidth}}{
        \cellcolor[rgb]{0.9,0.9,0.9}{
            \makecell[{{p{\linewidth}}}]{
                \texttt{\tiny{[GM$|$GM]}}
                \texttt{Error: The response should always start with the keyword 'explanation:'} \\
            }
        }
    }
    & & \\ \\

    \theutterance \stepcounter{utterance}  
    & & & \multicolumn{2}{p{0.3\linewidth}}{
        \cellcolor[rgb]{0.9,0.9,0.9}{
            \makecell[{{p{\linewidth}}}]{
                \texttt{\tiny{[GM$|$GM]}}
                \texttt{game\_result = ABORT} \\
            }
        }
    }
    & & \\ \\

\end{supertabular}
}

\end{document}
